$\dim(\mathcal{A}) + \dim(\mathcal{A}^{\perp}) = |U|$, for all subspaces $\mathcal{A} \subseteq \

PROOF. Since $\mathcal{A}$ is a subspace, $\mathcal{A} = (\mathcal{A}^{\perp})^{\perp}$. Hence:

$$U \in \mathcal{A} \Leftrightarrow S \cdot U = 0 \quad \text{for all } S \in \mathcal{A}^{\perp}$$
$$\Leftrightarrow |S \cap U| \equiv 0 \text{ (modulo 2)} \quad \text{for all } S \in \mathcal{A}^{\perp}$$
$$\Leftrightarrow |S| \quad \text{is even for all } S \in \mathcal{A}^{\perp}$$

If $\mathcal{S} \subseteq \mathcal{P}(U)$, the foundation of $\mathcal{S}$ is the set $\text{Fnd}(\mathcal{S}) = \bigcup_{s \in \mathcal{S}} S$. If $\mathcal{A}$ and $\mathcal{B}$ are subspace of $\mathcal{P}(U)$, then clearly their intersection $\mathcal{A} \cap \mathcal{B}$ is a subspace of $\mathcal{P}(U)$. Their join (also referred to as “sum”) given by

$$\mathcal{A} \vee \mathcal{B} = \{S + T: S \in \mathcal{A}; T \in \mathcal{B}\}$$

is also a subspace of $\mathcal{P}(U)$. Clearly $\vee$ is a commutative and associative operation on the collection of subspaces of $\mathcal{P}(U)$. In the special case where the foundations of $\mathcal{A}$ and $\mathcal{B}$ are disjoint, the subspace $\mathcal{A} \vee \mathcal{B}$ is called the coordinate sum of $\mathcal{A}$ and $\mathcal{B}$ and is denoted by $\mathcal{A} \oplus \mathcal{B}$.

The following is a standard result from linear algebra:
